% mazzi: 3   % le stesse slide sono le pp. 14--28 del mazzo 02, lasciate al cap. 02
\chapter{La possibile rilevanza economica degli impianti}

Stabilito che la gestione di un impianto sportivo è di regola una concessione di servizi, il nuovo
codice dei contratti pubblici impone una seconda distinzione: fra \textbf{servizi a rilevanza
economica} e \textbf{servizi privi di tale rilevanza}. Da essa dipendono le modalità e le
procedure concrete di affidamento.

% ---------------------------------------------------------------------------
\section{Il nuovo codice dei contratti pubblici}

\textbf{D.Lgs.\ 50/2016} (in vigore da aprile 2016): abroga il D.Lgs.\ 163/2006 e ne rovescia
l'impostazione. Il vecchio testo era di fatto un codice dei soli appalti, perché escludeva le
concessioni di servizi; il nuovo le include.

\textbf{Art.\ 1}: ``il presente codice disciplina i contratti di appalto \textbf{e di concessione}
delle amministrazioni aggiudicatrici e degli enti aggiudicatori aventi ad oggetto l'acquisizione di
servizi, forniture, lavori e opere, nonché i concorsi pubblici di progettazione''.

\begin{itemize}
  \item \textbf{conseguenza}: con l'inclusione delle concessioni di servizi fra le tipologie
    contrattuali disciplinate dal codice, nasce l'esigenza di verificare quale disciplina si
    applichi davvero all'affidamento in gestione degli impianti sportivi pubblici, prima
    espressamente esclusi.
\end{itemize}

\rubrica{Parte III --- contratti di concessione, art.\ 164}

\begin{itemize}
  \item \textbf{comma 2}: alle procedure di aggiudicazione di concessioni di lavori pubblici o di
    servizi si applicano, \textbf{per quanto compatibili}, le disposizioni delle parti I e II del
    codice --- principi generali, esclusioni, modalità e procedure di affidamento, pubblicazione e
    redazione di bandi e avvisi, requisiti generali e speciali e motivi di esclusione, criteri di
    aggiudicazione, comunicazioni a candidati e offerenti, requisiti di qualificazione degli
    operatori economici, termini di ricezione delle domande e delle offerte, modalità di
    esecuzione;
  \item \textbf{comma 3}: ``i \textbf{servizi non economici di interesse generale} non rientrano
    nell'ambito di applicazione della presente Parte'' $\rightarrow$ è la disposizione da cui
    dipende tutto il seguito, perché sottrae alla disciplina delle concessioni i servizi privi di
    rilevanza economica.
\end{itemize}

% ---------------------------------------------------------------------------
\section{La classificazione dei servizi pubblici}

I servizi pubblici si dividono in \textbf{servizi a rilevanza economica} e \textbf{servizi privi di
rilevanza economica}. La distinzione è legata:

\begin{itemize}
  \item all'\textbf{impatto} che l'attività può avere sull'assetto della \textbf{concorrenza};
  \item ai suoi caratteri di \textbf{redditività}.
\end{itemize}

\textbf{Limiti della classificazione}: è \textbf{solo astratta}, priva di carattere dinamico ed
evolutivo, e non può essere considerata assoluta rispetto a una considerazione globale
dell'ipotesi concreta, che va valutata:

\begin{enumerate}
  \item dal punto di vista del \textbf{servizio} che si affida (uno stadio per competizioni
    nazionali e internazionali, un campo di quartiere, un bocciodromo, una piscina, un impianto di
    tiro con l'arco non sono la stessa cosa);
  \item dal punto di vista del \textbf{soggetto destinatario} dell'affidamento, cioè
    dell'operatore economico --- anche quando sia privo di scopo di lucro;
  \item dal \textbf{metodo} con cui il servizio viene fornito.
\end{enumerate}

Non è quindi corretto chiedersi se la gestione di impianti sportivi abbia \emph{sempre} o
\emph{mai} rilevanza economica.

% ---------------------------------------------------------------------------
\section{I criteri della rilevanza economica}

\subsection{Il punto di partenza e la sua smentita}

\textbf{Regola generale}: i servizi pubblici come quelli sportivi sono \textbf{generalmente
considerati privi di rilevanza economica}, perché si ritiene debbano essere resi alla collettività
anche al di fuori di una logica di profitto d'impresa, e perché non esiste una situazione di
mercato appetibile per gli imprenditori --- la sola gestione non consente di trarre una notevole
fonte di remunerazione.

\textbf{Cons.\ Stato n.\ 5072/2006}: ``anche i servizi cosiddetti sociali, connotati da
significativo rilievo socio-assistenziale possono risultare compatibili con la nozione di servizio
di rilevanza economica se e nel momento in cui presentano, per le modalità di esplicazione, una
rilevante componente economica, tesa ad assicurare non solo la mera copertura delle spese sostenute
ma anche un \textbf{potenziale profitto di impresa}''.

\subsection{I due criteri}

\begin{itemize}
  \item \textbf{di rilevanza economica}: il servizio che si innesta in un settore per il quale
    esiste, \textbf{quantomeno in potenza}, una redditività; la gestione consente una
    \textbf{remunerazione dei fattori produttivi e del capitale} ed è quindi possibile una
    competizione sul mercato --- e ciò \textbf{ancorché} siano previste forme di intervento
    finanziario pubblico (il contributo di gestione e le tariffe calmierate non escludono, da
    soli, la rilevanza economica);
  \item \textbf{privo di rilevanza}: il servizio che, per sua natura o per le modalità di
    gestione, \textbf{non dà luogo ad alcuna competizione}, è irrilevante ai fini della
    concorrenza, e la cui gestione deve rispondere al solo principio di \textbf{buon andamento} e
    non anche al necessario equilibrio fra costi e ricavi.
\end{itemize}

(Parere Corte dei Conti n.\ 195/2009; TAR Puglia n.\ 1318/2006; Cons.\ Stato n.\ 5072/2006.)

L'assenza di rilevanza economica \textbf{non} autorizza però l'affidamento diretto: resta
necessaria una procedura ad evidenza pubblica che ponga tutti gli operatori sullo stesso piano;
sono le modalità di affidamento del servizio a rilevanza economica a essere più stringenti.

\subsection{La valutazione caso per caso}

\textbf{Corte dei Conti, sez.\ controllo Sardegna, delib.\ n.\ 9 del 18 giugno 2007}: ``non è
possibile individuare a priori, in maniera definita e statica, una categoria di servizi pubblici a
rilevanza economica''; la valutazione va effettuata \textbf{di volta in volta}, con riferimento al
singolo servizio da espletare, da parte dell'ente stesso.

\begin{itemize}
  \item \textbf{parametri}: l'impatto del servizio sullo specifico mercato concorrenziale di
    riferimento e i suoi caratteri di \textbf{redditività/autosufficienza economica}, cioè la
    capacità di produrre profitti o quantomeno di coprire i costi con i ricavi;
  \item \textbf{compito dell'ente}: valutare le modalità ottimali di espletamento del servizio
    quanto a costi, margini di copertura degli stessi e organizzazione in termini di efficienza,
    efficacia ed economicità, nel rispetto della tutela della concorrenza da un lato e
    dell'universalità e dei livelli essenziali delle prestazioni dall'altro;
  \item \textbf{conseguenza pratica}: i regolamenti comunali che elencano in anticipo quali
    impianti abbiano rilevanza economica fotografano la situazione al momento in cui sono redatti
    $\rightarrow$ l'analisi va tenuta dinamica, non trattata come acquisita.
\end{itemize}

\rubrica{Come la rilevanza può cambiare nel tempo}

\begin{itemize}
  \item \textbf{concorrenza sopravvenuta}: una piscina comunale unico impianto natatorio nel raggio
    di alcune decine di chilometri ha redditività potenzialmente elevata; la costruzione di una
    piscina nuova e più efficiente in un comune limitrofo può privarla di quella rilevanza. Lo
    stesso vale per gli stadi, quando si costruisce ex novo invece di riqualificare l'esistente;
  \item \textbf{modalità di utilizzo}: la stessa piscina destinata ad attività corsistica intensa
    --- molti utenti paganti in acqua --- è redditizia; destinata ad attività prettamente
    agonistica, con un atleta per corsia, diventa priva di rilevanza economica.
\end{itemize}

\subsection{Ciò che non rileva: la natura dell'affidatario}

È \textbf{irrilevante}, ai fini della connotazione dell'impianto e della conseguente modalità di
affidamento, la natura del futuro soggetto affidatario --- associazione o società sportiva
dilettantistica oppure ente commerciale a tutti gli effetti.

\begin{itemize}
  \item non è il destino dell'utile --- \textbf{reinvestito} nell'attività per le finalità sociali
    oppure \textbf{ridistribuito} fra i soci --- a determinare la potenziale redditività di un
    servizio pubblico, ``anche se ciò ha spesso tratto in inganno'';
  \item l'equivoco ricorrente è ritenere che l'affidamento a un soggetto senza scopo di lucro
    escluda in partenza la rilevanza economica: il divieto di distribuire gli utili riguarda
    l'affidatario, non la natura del servizio.
\end{itemize}

\subsection{Le variabili del settore sportivo}

Nel settore sportivo è impossibile operare distinzioni precise per tipologia di impianto, perché
troppe sono le variabili che incidono:

\begin{itemize}
  \item \textbf{grandezza}: le dimensioni dell'impianto;
  \item \textbf{bacino d'utenza}: non coincide con la popolazione del comune, perché l'utenza si
    sposta più facilmente dall'hinterland verso il capoluogo che nella direzione opposta
    $\rightarrow$ a parità di bacino potenziale, lo stesso impianto collocato in un piccolo comune
    vale meno nel piano economico;
  \item \textbf{modalità di gestione}.
\end{itemize}

\textbf{Studio caso per caso}: la possibile redditività va determinata analizzando tutti i fattori
capaci di indicarne la potenziale produttività --- tipologie di attività praticate, costi del
personale, tariffe praticabili, presenza di spazi per il pubblico (e quindi possibilità di
spettacolo sportivo o di intrattenimento) e di spazi per la cartellonistica pubblicitaria.

\subsection{La sintesi}

\begin{itemize}
  \item \textbf{impianti senza rilevanza economica}: la gestione \textbf{va assistita dall'ente
    pubblico}, perché non è in grado di sostenersi da sola;
  \item \textbf{impianti con rilevanza economica}: la gestione è in grado di \textbf{sostenersi e
    di produrre reddito}.
\end{itemize}

\textbf{TAR Marche sez.\ I n.\ 73/2013}: il servizio di gestione di impianti sportivi pubblici
(nella specie una piscina comunale) che produce utili economici anche rilevanti riveste natura di
servizio pubblico \textbf{a rilevanza economica}.

\begin{itemize}
  \item non basta a escluderla il fatto che le \textbf{tariffe} debbano essere concordate con il
    Comune, né l'esistenza di altri vincoli a carico dell'aggiudicatario: rispondono alla finalità
    ultima del servizio, cioè consentire a società sportive e singoli utenti di utilizzare
    l'impianto;
  \item la sottoposizione delle tariffe all'approvazione dell'autorità concedente è del resto
    ``un fenomeno del tutto comune in materia di concessioni di beni e servizi pubblici''.
\end{itemize}

% ---------------------------------------------------------------------------
\section{La disciplina degli impianti privi di rilevanza economica}

Quale disciplina applicare quando l'impianto è privo di rilevanza economica, cioè quando si tratta
di un \textbf{``servizio non economico di interesse generale''} escluso dalla parte III del codice
ex art.\ 164 comma 3? Le ipotesi sul tavolo sono \textbf{due}:

\begin{enumerate}
  \item si applica la disciplina prevista per gli \textbf{appalti di servizi};
  \item si applicano \textbf{unicamente} i principi comuni della sola parte I del codice e la
    parte IV relativa al PPP --- fra le cui forme è inserita la concessione di servizi --- e non
    anche le ulteriori previsioni della parte III, che disciplina nello specifico i contratti di
    concessione.
\end{enumerate}

\subsection{Prima ipotesi: la delibera ANAC 1300/2016}

\textbf{Delibera ANAC n.\ 1300 del 14 dicembre 2016}, resa su quesito in materia di affidamento di
servizi pubblici privi di rilevanza economica, ribalta la prassi fino ad allora consolidata.

\begin{enumerate}
  \item secondo il \textbf{vocabolario comune per gli appalti pubblici (CPV)} --- Reg.\ (CE)
    n.\ 2195/2002, come modificato dal Reg.\ (CE) n.\ 213/2008 --- il codice CPV
    \textbf{92610000-0} è riferito ai ``servizi di gestione di impianti sportivi'';
  \item tale CPV è ricompreso nell'\textbf{Allegato IX} del D.Lgs.\ 50/2016 (servizi di cui agli
    artt.\ 140, 143 e 144), nella categoria ``servizi amministrativi, sociali, in materia di
    istruzione, assistenza sanitaria e cultura'';
  \item si tratterebbe quindi di un \textbf{appalto di servizi}, perché oggetto dell'affidamento è
    la gestione dell'impianto quale servizio reso per conto dell'amministrazione e \textbf{in
    assenza di rischio operativo} (art.\ 3 del codice);
  \item ne discende che la gestione degli impianti privi di rilevanza economica, sottratta alla
    disciplina delle concessioni, va ricondotta agli ``appalti di servizi'' e aggiudicata secondo
    le previsioni dettate per gli \textbf{appalti di servizi sociali} (Titolo VI, sez.\ IV), ferma
    la disciplina dell'art.\ 36 per gli affidamenti sotto le soglie dell'art.\ 35;
  \item \textbf{corollario}: dovrebbe ritenersi \textbf{superata e non più applicabile} la
    previsione dell'art.\ 90 comma 25 L.\ 289/2002 --- l'affidamento in via preferenziale a
    società e associazioni sportive dilettantistiche --- perché dettata in un differente contesto
    normativo.
\end{enumerate}

\rubrica{Obiezioni}

\begin{itemize}
  \item il ragionamento presuppone che il servizio privo di rilevanza economica sia \textbf{privo
    di rischio operativo}, mentre il rischio di perdita gestionale resta in capo al gestore: non
    esistono più, da tempo, convenzioni che prevedano la copertura del disavanzo da parte
    dell'Amministrazione;
  \item ne deriverebbe il \textbf{paradosso} di regole più restrittive per l'affidamento dei
    servizi privi di rilevanza economica --- soggetti in tutto alle norme sull'appalto --- che per
    le concessioni a rilevanza economica;
  \item qualificare l'affidamento come appalto implicherebbe inoltre porre il \textbf{rischio
    economico interamente in capo alla PA}, ipotesi non rispondente alla realtà, dato che gli
    impianti sono sempre più affidati alla gestione esclusiva del privato.
\end{itemize}

\rubrica{L'esito: la norma non è abrogata, ma va letta restrittivamente}

L'art.\ 90 comma 25 continua a reggere le leggi regionali sull'affidamento (cap.~4). Il criterio di
preferenzialità si porrebbe in violazione dei principi comunitari e delle nuove norme interne in
materia di contrattualistica pubblica \textbf{solo laddove si traducesse in un semplicistico
affidamento diretto} senza alcuna forma di \textbf{pubblicità preventiva}, che resta dunque
necessaria in ogni caso --- \textbf{anche} quando l'impianto sia privo di rilevanza economica, dove
l'affidamento non può che avvenire nel rispetto dei principi desumibili dal Trattato.

\subsection{Seconda ipotesi: gli artt.\ 4 e 30 del codice}

È la lettura considerata l'unica praticabile fino a tutto il 2016. Poiché l'art.\ 164 comma 3
esclude i servizi non economici di interesse generale dalla parte III, essi sono contratti
\textbf{esclusi in parte} dall'ambito di applicazione del codice, e ricadono quindi nell'art.\ 4.

\textbf{Art.\ 4 --- principi relativi all'affidamento di contratti pubblici esclusi}:
l'affidamento dei contratti pubblici aventi ad oggetto lavori, servizi e forniture esclusi, in
tutto o in parte, dall'ambito di applicazione oggettiva del codice avviene nel rispetto dei
principi di:

\begin{multicols}{2}
\begin{itemize}
  \item economicità;
  \item efficacia;
  \item imparzialità;
  \item parità di trattamento;
  \item trasparenza;
  \item proporzionalità;
  \item pubblicità;
  \item tutela dell'ambiente;
  \item efficienza energetica.
\end{itemize}
\end{multicols}

Gli ultimi due sono l'aggiunta rispetto alla formulazione dell'art.\ 30 del vecchio codice.

\textbf{Art.\ 30 --- principi per l'aggiudicazione e l'esecuzione di appalti e concessioni}:
l'affidamento e l'esecuzione garantiscono la \textbf{qualità delle prestazioni} e si svolgono nel
rispetto dei principi di economicità, efficacia, tempestività e correttezza; le stazioni appaltanti
rispettano inoltre i principi di libera concorrenza, non discriminazione, trasparenza e
proporzionalità, oltre alla pubblicità.

\begin{itemize}
  \item \textbf{novità rilevante}: il principio di \textbf{economicità può essere subordinato} ---
    nei limiti in cui è espressamente consentito dalle norme vigenti e dal codice --- a criteri,
    previsti nel bando, ispirati a \textbf{esigenze sociali}, nonché alla tutela della salute,
    dell'ambiente e del patrimonio culturale e alla promozione dello sviluppo sostenibile, anche
    dal punto di vista energetico;
  \item è la disposizione che consente di non affidare il servizio secondo criteri puramente
    economici proprio là dove le esigenze sociali connotano l'offerta sportiva e la gestione di un
    impianto pubblico.
\end{itemize}
